\documentclass[11pt]{article}
\usepackage{laborstil}
\hypersetup{
  pdftitle={TP4: Spektrogramm und Sprachaktivität},
  pdfsubject={Labor Versuch TP4: Frequenzen über der Zeit und Sprachaktivitätserkennung},
  pdfkeywords={Digitale Audiosignalverarbeitung, Python, Labor}
}
\fusszeile{Digitale Audiosignalverarbeitung, Labor Versuch TP4}

\begin{document}

\titelseite{Spektrogramm und Sprachaktivität:\\ Frequenzen über der Zeit}{Labor Versuch TP4}{TP4\_Spektrogramm\_Sprachaktivitaet.ipynb}

% ---------------------------------------------------------------------
\section{Hinweise zur Durchführung}

Gerechnet, gezeichnet und angehört wird im zugehörigen Jupyter-Notebook (Dateiname auf der Titelseite). Dieses Heft führt durch den Versuch. Abschnitte und Aufgaben haben im Heft und im Notebook \textbf{dieselben Nummern}. Das Notebook braucht kein anderes Notebook.

\notebooklink{qr/TP4_Spektrogramm_Sprachaktivitaet_colab.pdf}{https://colab.research.google.com/github/Lamhour-Mohamed-Akram/audio-dsp-lab/blob/main/TP4_Spektrogramm_Sprachaktivitaet.ipynb}{https://colab.research.google.com/github/Lamhour-Mohamed-Akram/}{audio-dsp-lab/blob/main/TP4\_Spektrogramm\_Sprachaktivitaet.ipynb}

\begin{itemize}
  \item \textbf{Öffnen:} Link oder QR-Code öffnen, gegebenenfalls mit einem Google-Konto anmelden und die Zellen von oben nach unten mit \textbf{Shift + Enter} ausführen. Eine Warnung, dass das Notebook nicht von Google stammt, mit "`Trotzdem ausführen"' bestätigen.
  \item \textbf{Antworten:} in dieses Heft \textbf{oder} in eine eigene Kopie des Notebooks (\textit{Datei > Kopie in Drive speichern}). Schreibrechte für das Original sind nicht nötig.
  \item \textbf{Pflicht und Zusatz:} Pflichtaufgaben bearbeiten alle, Zusatzaufgaben sind freiwillig. \textbf{Zeitbedarf (Schätzung):} Pflichtaufgaben 4.1 bis 4.4 ca. 60 bis 90 Minuten nach der Einrichtung.
  \item \textbf{Arbeitsweise:} immer nur ein oder zwei Parameter ändern, Zelle erneut ausführen, hören, beobachten und notieren. Kopfhörer verwenden.
\end{itemize}

\subsection{Einrichtung}

Die ersten zwei Code-Zellen installieren nur fehlende Python-Pakete und laden \texttt{speech.wav}, \texttt{music.wav} und \texttt{sound.wav} automatisch in den Ordner \texttt{data/}, falls sie fehlen oder beschädigt sind. Uploads, Google Drive, ein GitHub-Konto oder eine GPU sind nicht nötig. \textbf{Neustart:} Nach \textit{Sitzung neu starten} sind die Variablen weg, die Dateien bleiben. Nach \textit{Laufzeit trennen und löschen} oder in einer neuen Laufzeit sind auch die Dateien weg. Dann Abschnitt 1 und alle Zellen bis zur aktuellen Stelle erneut ausführen, am einfachsten mit \textit{Laufzeit > Alle ausführen}.

\subsection{Bibliotheken und Hilfsfunktionen}

Eine Zelle lädt die Bibliotheken und die aus TP2 bekannten Hilfsfunktionen, zum Beispiel \texttt{anhoeren} (ohne automatische Lautstärke-Anpassung), \texttt{gemeinsame\_skala} und \texttt{spektrum}.

\section{Lernziele}

Nach diesem Versuch können Sie
\begin{itemize}
  \item den Unterschied zwischen dem Spektrum der ganzen Aufnahme und einem Spektrogramm erklären,
  \item Spektrogramme von Sprache, Musik und Geräusch lesen und vergleichen,
  \item ein Mel-Spektrogramm in einfachen Worten erklären,
  \item mit der Kurzzeit-Energie ungefähr erkennen, wann gesprochen wird,
  \item zeigen, wann diese Methode versagt (Rauschen, leise Sprache),
  \item Sprachaktivitätserkennung, Sprechererkennung und Spracherkennung unterscheiden.
\end{itemize}

% ---------------------------------------------------------------------
\bigskip
\section{Vorbereitungsaufgaben}

Bearbeiten Sie diese Aufgaben \textbf{vor} dem Labor. Kurze Antworten genügen.

\begin{block}
\textbf{V1:} Ein Spektrogramm benutzt Fenster von 32\,ms. Wie viele Samples sind das bei 16000\,Hz?
\antwortlinien{1}
\end{block}
\begin{block}
\textbf{V2:} Das Spektrum der ganzen Aufnahme zeigt nicht, wann eine Frequenz vorkommt. Warum ist das bei Sprache ein Nachteil?
\antwortlinien{2}
\end{block}
\begin{block}
\textbf{V3:} Menschen hören den Unterschied zwischen 100\,Hz und 200\,Hz sehr deutlich, zwischen 8000\,Hz und 8100\,Hz kaum. Was bedeutet das für eine gehörgerechte Frequenzachse?
\antwortlinien{2}
\end{block}
\begin{block}
\textbf{V4:} Erklären Sie den Unterschied zwischen den Fragen: Wann spricht jemand? Wer spricht? Was wird gesagt?
\antwortlinien{2}
\end{block}

% ---------------------------------------------------------------------
\bigskip
\section{Durchführung}

\subsection[Vom Spektrum zum Spektrogramm]{Vom Spektrum zum Spektrogramm (Aufgabe 4.1, Pflicht)}

\ziel{Sehen, wie sich der Frequenzinhalt über der Zeit ändert.}

Die FFT der ganzen Aufnahme mischt alle Zeitpunkte zusammen. Ein \textbf{Spektrogramm} zerlegt das Signal in kurze, überlappende Fenster, berechnet für jedes Fenster eine FFT und zeigt die Spektren nebeneinander: x-Achse Zeit, y-Achse Frequenz, Farbe Stärke (hell = stark, 0\,dB = stärkster Punkt des Bildes). Parameter: \param{n\_fft} ist die Fensterlänge in Samples (gültig 64 bis 8192), \param{hop} der Schritt von Fenster zu Fenster (gültig 1 bis \param{n\_fft}). Voreinstellung: \param{n\_fft = 512} (32\,ms bei 16000\,Hz) und \param{hop = 128}.

\begin{block}
\aufgabe{4.1}{Pflicht}
\begin{enumerate}
  \item Finden Sie im Spektrogramm die Pausen (dunkle senkrechte Streifen). Nennen Sie zwei Zeitpunkte.
  \item In Vokalen sieht man viele waagerechte Linien übereinander. Was stellen sie dar?
  \item Setzen Sie \param{n\_fft = 128} und danach \param{n\_fft = 2048}. Was wird jeweils schärfer: die zeitliche Darstellung oder die Frequenzdarstellung?
\end{enumerate}
\antwortlinien{2}
\end{block}

\subsection[Sprache, Musik, Geräusch und Mel-Spektrogramm]{Sprache, Musik, Geräusch und Mel-Spektrogramm (Aufgabe 4.2, Pflicht)}

\ziel{Typische Muster erkennen und die Mel-Skala kennenlernen.}

\begin{itemize}
  \item Das Notebook zeigt die Spektrogramme aller drei Aufnahmen untereinander. Stereo wird nur für die Darstellung zu Mono gemittelt.
  \item Das Gehör unterscheidet tiefe Frequenzen fein und hohe grob. Die \textbf{Mel-Skala} bildet das nach. Ein \textbf{Mel-Spektrogramm} fasst die FFT-Frequenzen zu wenigen Mel-Bändern zusammen. Solche Bilder werden häufig als Eingabe für KI-Modelle benutzt, zum Beispiel für neuronale Netze, die Audio wie ein Bild verarbeiten.
  \item Parameter: \param{n\_mels}, die Anzahl der Mel-Bänder (gültig 8 bis 128, Voreinstellung 64).
\end{itemize}

\begin{block}
\aufgabe{4.2}{Pflicht}
\begin{enumerate}
  \item Beschreiben Sie in \textbf{Tabelle A} das typische Muster jeder Aufnahme. Achten Sie auf Pausen, lange waagerechte Linien (gehaltene Töne), senkrechte Striche (Schläge, Knacken), steigende oder fallende Linien (Tonhöhe ändert sich) und breites Rauschen.
  \item Setzen Sie \param{n\_mels = 20} und danach \param{n\_mels = 128}. Wie verändert sich das Mel-Spektrogramm? Sind Pausen und Vokale noch zu erkennen?
\end{enumerate}
\antwortlinien{2}
\end{block}

\subsection[Sprachaktivität mit der Kurzzeit-Energie]{Sprachaktivität mit der Kurzzeit-Energie (Aufgabe 4.3, Pflicht)}

\ziel{Ungefähr erkennen, wann in \texttt{speech.wav} gesprochen wird (Voice Activity Detection, VAD).}

\begin{enumerate}[label=\arabic*.]
  \item Das Signal wird in überlappende Fenster zerlegt (25\,ms lang, Schritt 10\,ms). Das letzte Fenster endet genau am Ende der Aufnahme.
  \item Für jedes Fenster wird die Energie berechnet: der Mittelwert der quadrierten Samples.
  \item Die Energie wird in dB umgerechnet, 0\,dB ist das lauteste Fenster. Fenster ohne jede Energie erhalten minus unendlich dB.
  \item Fenster über der Schwelle \param{schwelle\_db} gelten als aktiv (gültig -100 bis 0\,dB).
  \item Aktive Fenster, die sich überlappen oder berühren, bilden einen Abschnitt vom Anfang des ersten bis zum Ende des letzten aktiven Fensters.
\end{enumerate}
Die Grenzen sind \textbf{ungefähr} und entsprechen nicht genau den Wortgrenzen. Findet die Methode nichts, meldet das Notebook das, statt abzubrechen.

\begin{block}
\aufgabe{4.3}{Pflicht}
Probieren Sie \param{schwelle\_db} = -10, -20, -30, -40 und -50. Tragen Sie in \textbf{Tabelle B} die Anzahl der Abschnitte, den Anteil aktiver Fenster und Ihre Beobachtung ein: Fehlen Teile von Wörtern? Werden Pausen als aktiv markiert? Hören Sie dazu auch die aneinandergehängten aktiven Abschnitte an.
\end{block}

\subsection[Grenzen der Methode]{Grenzen der Methode (Aufgabe 4.4, Pflicht)}

\ziel{Zeigen, wann eine Energie-Schwelle versagt.}

\begin{itemize}
  \item \textbf{a) Rauschen:} Das Notebook erzeugt die verrauschte Sprache selbst (\param{seed=42}, \param{vad\_snr\_db = 10}, Vorschläge 20, 10, 0) und wendet \param{schwelle\_rauschen\_db = -30} an.
  \item \textbf{b) Leise Sprache:} Die zweite Hälfte wird mit \param{leise\_faktor = 0.03} (etwa -30\,dB, gültig 0 bis 1) leiser gemacht. Die Schwelle \param{schwelle\_leise\_db = -30} bezieht sich auf das lauteste Fenster der ganzen Aufnahme.
\end{itemize}

\begin{block}
\aufgabe{4.4}{Pflicht}
\begin{enumerate}
  \item Beobachten Sie bei \param{vad\_snr\_db = 10} und \param{schwelle\_rauschen\_db = -30}, welche Bereiche als aktiv markiert werden. Suchen Sie eine Schwelle, bei der die Pausen wieder als inaktiv erkannt werden (\textbf{Tabelle C}).
  \item Welche Teile der leisen zweiten Hälfte werden bei \param{schwelle\_leise\_db = -30} erkannt? Suchen Sie eine Schwelle, bei der die leise Sprache erkannt wird.
  \item Gibt es eine gemeinsame Schwelle, die für beide Aufnahmen gut funktioniert? Was folgt daraus für diese einfache Methode?
\end{enumerate}
\antwortlinien{2}
\end{block}

\subsection{Wann, wer, was? Drei verschiedene Aufgaben}

{\small
\begin{tabularx}{\linewidth}{|L{3.9cm}|L{2.4cm}|Y|L{3.0cm}|}
  \hline
  \textbf{Aufgabe} & \textbf{Frage} & \textbf{Beispiel für ein Ergebnis} & \textbf{In diesem Labor} \\ \hline
  Sprach\-aktivitäts\-erkennung (VAD) & Wann wird gesprochen? & "`aktiv von 0,4\,s bis 2,1\,s"' & ja, einfach, mit der Kurzzeit-Energie \\ \hline
  Sprechererkennung (speaker identification) & Wer spricht? & "`Person A"' & nein \\ \hline
  Spracherkennung (ASR) & Was wird gesagt? & der gesprochene Text & nein \\ \hline
\end{tabularx}}

\medskip
Eine Energie-Schwelle unterscheidet nur zwischen viel und wenig Energie: Laute Geräusche können als Sprache gelten, leise Sprache kann übersehen werden. Sie sagt nichts darüber, wer spricht oder was gesagt wird. \texttt{speech.wav} enthält eine einzige Sprecherstimme, damit kann man Sprechererkennung weder üben noch testen. Beide brauchen trainierte Modelle, in diesem Labor wird kein Modell trainiert.

\subsection[Zusatzaufgabe Z4.1: Echo]{Zusatzaufgabe Z4.1: Echo (optional)}

Ein Echo ist eine verzögerte, leisere Kopie: Ausgang = Original + \param{daempfung} mal (Original, um \param{verzoegerung\_ms} später). Gültig: \param{verzoegerung\_ms} 1 bis 2000, \param{daempfung} 0 bis 1. Vergleichen Sie \param{verzoegerung\_ms} = 50, 300 und 1000 sowie \param{daempfung} = 0.3 und 0.8. Wann hören Sie ein deutliches Echo? Wann klingt es eher wie ein Raum oder eine Klangfärbung?

\subsection[Zusatzaufgabe Z4.2: MFCC]{Zusatzaufgabe Z4.2: MFCC (optional)}

MFCCs fassen jedes Fenster des Mel-Spektrogramms in wenigen Zahlen zusammen (\param{n\_mfcc = 13}, gültig 2 bis 40) und beschreiben grob die Klangfarbe. Sie sind klassische Eingaben für Verfahren des Maschinellen Lernens. Vergleichen Sie das MFCC-Bild mit dem Mel-Spektrogramm aus 4.2. Wo erkennen Sie die Pausen?

% ---------------------------------------------------------------------
\clearpage
\section{Auswertung und Diskussion}

\tabtitel{Tabelle A: Muster im Spektrogramm (Aufgabe 4.2 a)}
\begin{tabularx}{\linewidth}{|L{2.4cm}|Y|}
  \hline
  \textbf{Aufnahme} & \textbf{Typisches Muster} \\ \hline
  Sprache\rule{0pt}{1.1cm} & \\ \hline
  Musik\rule{0pt}{1.1cm} & \\ \hline
  Rennwagen\rule{0pt}{1.1cm} & \\ \hline
\end{tabularx}

\tabtitel{Tabelle B: Schwelle der Sprachaktivitätserkennung für speech.wav (Aufgabe 4.3)}
{\small
\begin{tabularx}{\linewidth}{|L{2.1cm}|L{2.0cm}|L{2.4cm}|Y|}
  \hline
  \textbf{schwelle\_\allowbreak db} & \textbf{Anzahl Abschnitte} & \textbf{Anteil aktiver Fenster [\%]} & \textbf{Beobachtung (Wortteile fehlen? Pausen aktiv?)} \\ \hline
  -10\schreibzeile & & & \\ \hline
  -20\schreibzeile & & & \\ \hline
  -30\schreibzeile & & & \\ \hline
  -40\schreibzeile & & & \\ \hline
  -50\schreibzeile & & & \\ \hline
\end{tabularx}}

\tabtitel{Tabelle C: Grenzen der Methode (Aufgabe 4.4)}
{\small
\begin{tabularx}{\linewidth}{|L{3.2cm}|L{3.2cm}|L{2.2cm}|Y|}
  \hline
  \textbf{Aufnahme} & \textbf{Einstellung} & \textbf{geeignete Schwelle [dB]} & \textbf{Beobachtung} \\ \hline
  verrauschte Sprache\schreibzeile & vad\_snr\_db = 10 & & \\ \hline
  zweite Hälfte leise\schreibzeile & leise\_faktor = 0.03 & & \\ \hline
\end{tabularx}}

\subsection*{Fragen}
\begin{enumerate}[label=\arabic*.]
  \frage{Was zeigt ein Spektrogramm, was das FFT-Spektrum der ganzen Aufnahme nicht zeigt?}{2}
  \frage{Was haben Sie bei \param{n\_fft = 128} und \param{n\_fft = 2048} beobachtet? Erklären Sie Ihre Beobachtung mit Zeit und Frequenz.}{2}
  \frage{Warum passt die Mel-Skala besser zum menschlichen Hören als eine lineare Frequenzachse?}{2}
  \frage{Welche Schwelle fanden Sie für \texttt{speech.wav} am geeignetsten? Wie hat sich dieselbe Schwelle bei der verrauschten und bei der leisen Aufnahme verhalten?}{2}
  \frage{Warum kann eine Energie-Schwelle nicht erkennen, wer spricht oder was gesagt wird?}{2}
  \frage{Nennen Sie eine Anwendung, in der Sprachaktivitätserkennung nützlich ist.}{1}
\end{enumerate}

% ---------------------------------------------------------------------
\section{Quellen und Audiodaten}

Die Aufnahmen stammen von Dritten. Für dieses Labor wurden nur Ausschnitte herausgeschnitten und als WAV-Dateien gespeichert (20\,ms Ein- und Ausblenden an den Schnittstellen, sonst keine Bearbeitung). Alle Details stehen in \texttt{DATA\_SOURCES.md} im Repository \url{https://github.com/Lamhour-Mohamed-Akram/audio-dsp-lab}. Die Übersicht aller Versuche mit Links steht dort in der Datei \texttt{README.md}.

{\small
\begin{itemize}
  \item \textbf{speech.wav} (16000\,Hz, mono, 15,9\,s): LibriVox-Hörbuch "`Fabeln"' von Magnus Gottfried Lichtwer, Abschnitt "`Das Wiesel und die Hühner"', gelesen von marham63. Lizenz: Public Domain. Quelle: \url{https://archive.org/details/fabeln_1905_librivox}. Ausschnitt 12,8\,s bis 28,7\,s, umgetastet auf 16000\,Hz.
  \item \textbf{music.wav} (44100\,Hz, stereo, 20,0\,s): Julius Fučík, "`Entrance of the Gladiators"', gespielt von der United States Marine Band. Lizenz: Public Domain (Werk der US-Regierung). Quelle: Wikimedia Commons, Datei \href{https://commons.wikimedia.org/wiki/File:Julius_Fu%C4%8D%C3%ADk%27s_%22Entrance_of_the_Gladiators%22,_performed_by_the_United_States_Marine_Band.flac}{\textit{Julius Fučík's Entrance of the Gladiators, performed by the United States Marine Band.flac}} (\url{https://commons.wikimedia.org}). Ausschnitt 24\,s bis 44\,s.
  \item \textbf{sound.wav} (22050\,Hz, stereo, 8,0\,s): "`Audi R8 (2000)"', aufgenommen von Edvvc beim Goodwood Festival of Speed 2009. Lizenz: CC BY-SA 3.0, \url{https://creativecommons.org/licenses/by-sa/3.0/}. Quelle: \url{https://commons.wikimedia.org/wiki/File:Audi_R8_(2000).ogg}. Ausschnitt 4\,s bis 12\,s, umgetastet auf 22050\,Hz. Der Ausschnitt steht ebenfalls unter CC BY-SA 3.0.
\end{itemize}}

\end{document}
